\section{Primera vía de escalamiento: Scaling Laws del preentrenamiento}

\subsection{Tres ejes que escalan en conjunto}

La regularidad empírica de la etapa de preentrenamiento es la siguiente: bajo configuraciones razonables de datos y optimización, aumentar la \textbf{cantidad de cómputo de entrenamiento, el volumen de datos de entrenamiento y el número de parámetros del modelo} suele reducir la pérdida de prueba. La forma abstracta más usada es una ley de potencia:

$$
L(x)=L_{\infty}+A x^{-\alpha},\qquad \alpha>0.
$$

\begin{itemize}
    \item $L(x)$: la pérdida de prueba dada una escala de recursos $x$;
    \item $L_{\infty}$: la cota inferior de pérdida difícil de eliminar bajo esa distribución de datos y esos supuestos de modelado;
    \item $A$: una constante de proporcionalidad relacionada con la tarea, los datos y la familia de modelos;
    \item $x$: puede representar cómputo, volumen de datos o número de parámetros;
    \item $\alpha$: el exponente de la ley de potencia, que determina la velocidad a la que baja la pérdida al aumentar los recursos.
\end{itemize}

La importancia de la ley de potencia no radica en que “siempre sea válida”, sino en que en su momento permitió a los equipos predecir, antes de entrenar, el beneficio de un modelo más grande, y con base en eso invertir en cómputo e ingeniería de datos a mayor escala.

\begin{figure}[H]
\centering
\includegraphics[width=0.96\textwidth]{figures/scaling-laws.jpg}
\caption{Scaling laws del preentrenamiento: los tres ejes de escalamiento de cómputo, datos y parámetros.\protect\footnotemark}
\end{figure}
\footnotetext{Intervalo del video: 00:02:50--00:03:53.}

\begin{knowledgebox}{“Más grande” no solo significa más parámetros}
Un modelo con un número de parámetros extremadamente grande, si carece de datos suficientes, terminará memorizando repetidamente un conjunto limitado de muestras; si el cómputo de entrenamiento es insuficiente, los parámetros no se optimizarán del todo; y si la calidad de los datos es mala, el cómputo adicional solo servirá para ajustar el ruido con mayor precisión. Por eso, el núcleo del scaling moderno es la \textbf{configuración coordinada de parámetros, tokens y cómputo efectivo}.
\end{knowledgebox}

\subsection{De BERT y GPT-2 a GPT-4: la escala como fuente de capacidad}

Entre 2,018 y 2,024, la escala de los modelos de lenguaje dominantes creció con rapidez. El curso usa esta historia para subrayar dos puntos: primero, que aumentar la escala trajo consigo una mejora continua del desempeño promedio; segundo, que ciertas capacidades no se vuelven visibles de forma gradual en modelos pequeños, sino que se convierten de golpe en un comportamiento observable al alcanzar cierta escala.

\begin{figure}[H]
\centering
\includegraphics[width=0.92\textwidth]{figures/model-growth.jpg}
\caption{El rápido crecimiento de la escala de parámetros de los modelos de lenguaje grandes.\protect\footnotemark}
\end{figure}
\footnotetext{Intervalo del video: 00:03:55--00:04:46.}

Aquí conviene evitar un malentendido común: que la escala de parámetros crezca en la figura no significa que los parámetros sean la única variable causal. La proporción de los datos, la arquitectura, el optimizador, la estabilidad del entrenamiento y los métodos de post-entrenamiento cambiaron todos en el mismo periodo, de modo que la curva histórica expresa el \textbf{resultado del escalamiento integral de un sistema de ingeniería}.

\subsection{Zero-shot, few-shot y aprendizaje en contexto}

Después de GPT-3, un cambio clave es que el modelo puede entender la tarea directamente a partir del prompt:

\begin{itemize}
    \item \textbf{Zero-shot}: solo se da la descripción de la tarea y la entrada, sin ejemplos;
    \item \textbf{One-shot}: además se da un ejemplo de entrada--salida;
    \item \textbf{Few-shot}: se dan unos cuantos ejemplos para que el modelo infiera del contexto el mapeo de la tarea.
\end{itemize}

Este tipo de capacidad se llama \textbf{in-context learning}: los pesos del modelo no se actualizan, pero el patrón presente en el contexto cambia el comportamiento de ese cómputo hacia adelante en particular.

\begin{figure}[H]
\centering
\includegraphics[width=0.94\textwidth]{figures/zero-few-shot.jpg}
\caption{Diferencia estructural entre los prompts de zero-shot y few-shot.\protect\footnotemark}
\end{figure}
\footnotetext{Intervalo del video: 00:05:40--00:06:38.}

\begin{warningbox}{El aprendizaje en contexto no es aprendizaje de parámetros}
Los ejemplos de few-shot no se escriben de forma permanente en los pesos durante la inferencia. Es más bien como establecer una interpretación temporal de la tarea para la solicitud actual. Una vez que termina el contexto, ese “resultado de aprendizaje” por lo general no se conserva de una solicitud a otra, a menos que el sistema cuente con memoria externa o un mecanismo de entrenamiento continuo.
\end{warningbox}

\subsection{Chain-of-Thought: llevar también el proceso intermedio al contexto}

Lo esencial del prompting de chain-of-thought (CoT) no es escribir unas cuantas frases más, sino aportar una \textbf{trayectoria de razonamiento intermedio que va del problema a la respuesta}. En comparación con mostrar únicamente la respuesta final, el modelo puede imitar la forma de descomponer el problema, actualizar el estado y realizar cálculos locales.

\begin{figure}[H]
\centering
\includegraphics[width=0.96\textwidth]{figures/chain-of-thought.jpg}
\caption{Comparación entre el prompting estándar y el prompting de chain-of-thought.\protect\footnotemark}
\end{figure}
\footnotetext{Intervalo del video: 00:07:18--00:09:01.}

El curso subraya que el efecto del CoT se refuerza claramente al aumentar la escala del modelo. En modelos pequeños, dar ejemplos de razonamiento puede aumentar solo la carga de imitación; en modelos suficientemente grandes, puede desencadenar un cómputo de varios pasos más confiable.

\begin{figure}[H]
\centering
\includegraphics[width=0.92\textwidth]{figures/cot-emergence.jpg}
\caption{La capacidad de chain-of-thought presenta un salto notable en modelos más grandes.\protect\footnotemark}
\end{figure}
\footnotetext{Intervalo del video: 00:09:01--00:10:18.}

\begin{importantbox}{De “predecir el siguiente token” a “ejecutar un proceso de cómputo”}
El CoT muestra que el mismo next-token predictor puede comportarse, dentro del contexto, como un calculador iterativo. El modelo no cambia su función objetivo, pero al generar tokens intermedios, descompone un mapeo difícil en varios mapeos locales más sencillos. Los reasoning model y los Agent posteriores retoman esta misma idea: \textbf{permitir que el sistema consuma más pasos intermedios antes de dar la respuesta final.}
\end{importantbox}

\subsection{Resumen de la sección}

El scaling del preentrenamiento amplía a la vez la capacidad promedio del modelo y el comportamiento potencial que puede activarse mediante el prompt. Zero-shot, few-shot y CoT muestran cómo el contexto puede reorganizar el cómputo del modelo sin actualizar los pesos, pero que el modelo “sepa algo” y que “actúe conforme a la intención humana” siguen siendo dos cosas distintas.
